\documentclass{article}

\textwidth=6.5in
\textheight=8.5in
\voffset=-54pt
\oddsidemargin=0in
\evensidemargin=0in
\setlength{\parskip}{8pt}
\setlength{\parindent}{0pt}

\usepackage{amsmath, amssymb}
\usepackage{ifthen}

\usepackage[inline]{enumitem}
\setlist{topsep=0pt, parsep=8pt}

\usepackage[rgb,table]{xcolor}\convertcolorsUfalse
\usepackage{graphicx}

% change hyperref options
\PassOptionsToPackage{hyphens}{url}
\usepackage[bookmarks,colorlinks,linkcolor=black,citecolor=blue,pdfstartview=FitH,urlcolor=blue]{hyperref}
\hypersetup{pdfpagemode=UseNone}
\usepackage{bookmark}

\usepackage{graphicx}
\usepackage{multicol}
\usepackage{framed}
\usepackage{array}

\usepackage{icomma}

\usepackage{soul}
\usepackage[normalem]{ulem}

\usepackage{pgfplots}
\pgfplotsset{compat=1.18}  
\pgfplotscreateplotcyclelist{mycolor}{black, blue, red, green!70!black, magenta, purple, cyan, orange }  
\pgfplotsset{
  disabledatascaling,
  cycle list=tikzcolor, 
  axis lines=middle,
  every axis plot/.style={no marks, thick},
  every axis/.style={
    enlarge x limits=false,
    enlarge y limits=auto,
    xlabel style={at={(current axis.right of origin)},anchor=west},
    ylabel style={at={(current axis.above origin)},anchor=south},
    % extra x and y ticks for asymptotes
    extra x tick style={grid=major, grid style={dashed,black}},
    extra y tick style={grid=major, grid style={dashed,black}},
    /pgf/number format/1000 sep={},
  },    
  tick label style = {font=\footnotesize, fill=\currentbgcolor, inner sep=0pt},    
  xticklabel shift=1pt,    
  scaled ticks=false,
  scaled y ticks=false,
  unbounded coords=jump,
  samples=101,
  minor tick length=0,
  scatter/classes={
    open={mark=*,draw=black,fill=white, mark size=2, thin},
    closed={mark=*,draw=black,fill=black, mark size=2, thin}
  },
  width=3in,
}
\pgfplotsset{open/.style={only marks, mark=*, fill=white, thin, mark size=2}}
\pgfplotsset{closed/.style={only marks, mark=*, draw=black, fill=black,
    thin, mark size=2}}
\pgfplotsset{labeled/.style={nodes near coords, only marks, mark=*, draw=black, fill=black, thin, mark size=2, point meta=explicit symbolic}}

\usepackage{tikz}
\usetikzlibrary{
  matrix,%
  % enables the snake paths
  decorations.pathmorphing,%
  % enables bigger arrows
  decorations.markings,%
  decorations.pathreplacing,%
  decorations.text,%
  calc,%
  patterns,%
  shapes.geometric,%
  arrows,
  decorations.shapes,
  positioning,
  plotmarks,
  intersections
}
\tikzset{>=latex} % sets default arrow type in tikz
\tikzset{font=\normalsize}
\tikzset{mark size=1.5pt, mark options=thin}
\tikzset{pin distance=4pt,
  every pin edge/.style={<-, thin, shorten <= -2pt}}

\interfootnotelinepenalty=10000
\renewcommand{\thefootnote}{(\arabic{footnote})}

\usepackage{multirow, bigdelim}

% change to \usepackage[sols]{handout} to see solutions
\usepackage[sols]{handout}

\newcommand\iscurrentenumi[1]{%
  \ifthenelse{\equal{#1}{\theenumi}}{}{#1}}
\setenumerate[1]{ref=\#\arabic*}
\setenumerate[2]{ref=\protect\iscurrentenumi{\theenumi}(\alph*)}
\newcommand\enumiiprefix[2]{%
  \ifthenelse{{\equal{#1}{\theenumi}}\AND{\equal{#2}{(\alph{enumii})}}}{}{}%
  \ifthenelse{{\equal{#1}{\theenumi}}\AND{\NOT{\equal{#2}{(\alph{enumii})}}}}{#2}{}%
  \ifthenelse{\equal{#1}{\theenumi}}{}{#1#2}%
}
\setenumerate[3]{ref=\protect\enumiiprefix{\theenumi}{(\alph{enumii})}\roman*}

\newlength{\tipmargin}
\makeatletter

\renewenvironment{framed}{
  \setlength{\tipmargin}{\textwidth-\linewidth}
  \advance\@totalleftmargin-\tipmargin
  \def\FrameCommand{\hspace{\tipmargin}\fbox}%
  \advance\linewidth\tipmargin
  \MakeFramed {\advance\hsize-\width \FrameRestore}%
}
{\endMakeFramed}

\makeatother

\definecolor{purple}{rgb}{0.5,0,1.0}

\usepackage{fancyhdr}
\lhead{\textcolor{white}{This content is protected and may not be shared, uploaded, or distributed.}}
\rhead{}
\rfoot{\vspace*{\fill}\footnotesize\textcopyright{} \emph{Harvard Math Ma}}
\renewcommand{\headrulewidth}{0pt}
\pagestyle{fancy}
\begin{document}

\htitle{35. Cubics and Higher-Degree Polynomials}

\begin{problemonly}
  \begin{framed}

You should be able to:
    \begin{itemize}[topsep=0pt,partopsep=0pt,parsep=8pt,itemsep=0pt]
    \item Identify the degree of a polynomial, and explain what this tells us about the possible number of roots and turning points of the polynomial. 

    \item Describe the ``end behavior'' of a polynomial, and explain what this tells us about the possible number of roots and turning points of the polynomial.

    \item Write a formula for a polynomial if given enough information about its characteristics (for example, roots, local extrema, etc.).

    \end{itemize}
  \end{framed}
\end{problemonly}

A \uline{polynomial} is any function that can be expressed in the form $f(x) = a_0 + a_1 x + a_2 x^2 + \cdots + a_n x^n$ where $a_0, \ldots, a_n$ are constants and the exponents are all non-negative integers. The \uline{degree} of a polynomial is the highest exponent in the polynomial.

\begin{enumerate}
\item  \textbf{Polynomial or not?}  Decide whether each of the following is
  a polynomial.  If so, what is its degree?
  \probonly{\begin{multicols}{2}}
    \begin{enumerate}
    \item
      \begin{problem}
        $f(x) = x^2 + 18x + 3$
      \end{problem}

      \begin{solution}
         Degree 2 polynomial.
      \end{solution}

    \item
      \begin{problem}
        $f(x) = x^{1/3} + \sqrt{x} + 5$
      \end{problem}

      \begin{solution}
        Not a polynomial.
      \end{solution}

    \item
      \begin{problem}
        $f(x) = e^x$
      \end{problem}

      \begin{solution}
        Not a polynomial.
      \end{solution}

    \item
      \begin{problem}
        $g(t) = -e t^4 + 8t - \pi$
      \end{problem}

      \begin{solution}
        Degree 4 polynomial.
      \end{solution}

    \item
      \begin{problem}
        $h(w) = 5w^3 + 7w + 9 w^{-1}$
      \end{problem}

      \begin{solution}
        Not a polynomial.
      \end{solution}

    \item
      \begin{problem}
        $j(u) = \ln \pi$
      \end{problem}

      \begin{solution}
        Degree 0 polynomial.
      \end{solution}
    \end{enumerate}

  \probonly{\end{multicols}}

\end{enumerate}
\begin{framed}
  As you know, a degree 2 polynomial is also called a \emph{quadratic}. A degree 3 polynomial is also called a \emph{cubic}.
\end{framed}

\begin{enumerate}[resume, topsep=0pt]
\item  \label{graph-cubic-with-repeated-root-followup}

  \note{Hopefully, students will find it straightforward to say that, if $g$ has 3 roots, then it should factor as $g(x) = (x - r_1)(x - r_2)(x - r_3)$.

    The case where $g$ has exactly one root is less obvious. I'll describe it as something like $(x - r_1) (\text{a quadratic with no real roots})$.

    Here, students should see graphically that it isn't possible for $g$ to have 0 roots. 
  }

  \begin{problem}[\newpage]
    In \href{https://people.math.harvard.edu/~jjchen/mathma/docs/homework/PS34.html}{Problem Set 34}, {{\#4}}, you found that the graph of $f(x) = (x-1)(x+3)^2 = x^3 + 5x^2 + 3x - 9$ looks like the picture below. If $d$ is a constant, how does the graph of $g(x) = x^3 + 5x^2 + 3x + d$ relate to the graph of $f$? How many roots could $g$ have?

    For each possible number of roots, sketch a possible graph of $g$, and say how you think the formula for $g(x)$ would factor.

    \begin{tikzpicture}
      \begin{axis}[xtick={1, -3}, ytick={-9}]
        \addplot+[domain=-5:2]{(x-1)*(x+3)^2};
      \end{axis}
    \end{tikzpicture}

  \end{problem}

  \begin{solution}
    The graph of $g(x)$ would be a vertical shift of the graph of
    $f(x)$ and could have 1, 2, or 3 real roots. If $g(x)$ has 3
    real roots, then it factors as $g(x)=(x-r_1)(x-r_2)(x-r_3)$. If
    $g(x)$ has two real roots, then it factors as $g(x)=(x-r_1)(x-r_2)^2$.
    If $g(x)$ has only one real root, then it factors as $g(x)=(x-r_1)(ax^2+bx+c)$, where $ax^2+bx+c$ has no real roots. Here are some
    possible graphs.
    \begin{center}
      \begin{tikzpicture}
        \begin{axis}[xtick={1, -3}, ytick=\empty]
          \addplot[domain=-5:2,violet]{(x-1)*(x+3)^2-5};
          \addplot[domain=-5:2]{(x-1)*(x+3)^2};
          \addplot[domain=-5:2,blue]{(x-1)*(x+3)^2+5};
          \addplot[domain=-5:2,red]{(x-1)*(x+3)^2+9.481};
          \addplot[domain=-5:2,orange]{(x-1)*(x+3)^2+15};
        \end{axis}
      \end{tikzpicture}
    \end{center}
    
  \end{solution}

\item  \emph{Comparing quadratics and cubics.}

  \begin{enumerate}
  \item

    \note{Ideally, students will have realized in the previous problem that the reason $g$ couldn't have 0 roots was that one end goes up while the other end goes down. Now, we'll generalize that.

      For a quadratic, I want them to see that $\displaystyle \lim_{x \to \infty} f(x) = \pm \infty$, and $\displaystyle \lim_{x \to -\infty} f(x) = \lim_{x \to \infty} f(x)$.

      For a cubic, we again have $\displaystyle \lim_{x \to \infty} f(x) = \pm \infty$, but now $\displaystyle \lim_{x \to -\infty} f(x) = - \lim_{x \to \infty} f(x)$.
    }

    \begin{problem}[\vfill]
      If $f(x)$ is a quadratic, what can you say about $\displaystyle \lim_{x \to \infty} f(x)$ and $\displaystyle \lim_{x \to -\infty} f(x)$? (We call this the ``end behavior'' of $f$.) What if $f(x)$ is a cubic?
    \end{problem}

    \begin{solution}
      If $f(x)$ is a quadratic, then its graph is a parabola, so
      $\lim\limits_{x\to-\infty} f(x) = \lim\limits_{x \to \infty} f(x)=
      \pm \infty$. If $f(x)$ is a cubic, then the ``ends'' go in
      opposite directions, so $\lim\limits_{x\to\-\infty} f(x)
      = -\lim\limits_{x\to\infty} f(x)$, so $\lim\limits_{x\to\infty} f(x)
      = \pm \infty$ and $\lim\limits_{x\to-\infty} f(x) = \mp \infty$.
    \end{solution}

  \item

    \note{We've already seen a cubic can have 1, 2, or 3 roots. The main question here is whether a cubic can have 0 roots, and the end behavior tells us that isn't possible.}

    \begin{problem}[\vfill]
      How many roots ($x$-intercepts) can a quadratic have? (Give all the possibilities.) How many roots can a cubic have? 
    \end{problem}

    \begin{solution}
      A quadratic can have 0, 1, or 2 real roots. A cubic can have
      1, 2, or 3 real roots. A cubic cannot have 0 real roots. Due to
      its end behavior, a cubic must cross the $x$-axis at least once.
    \end{solution}

  \item

    \note{The situation for cubics is a little complicated. A turning point of a cubic $f$ must have $f'(x) = 0$, and $f'$ is a quadratic. This could have 0, 1, or 2 roots, so we might start out thinking a cubic could have 0, 1, or 2 turning points. But if it had 1 turning point, the end behavior wouldn't be right, so it can only have 0 or 2 turning points. I will see if they can think of a cubic with no turning points ($x^3$!).}

    \begin{problem}[\vfill]
      How many turning points can a quadratic have? (Again, give all the possibilities.) How many turning points can a cubic have? 
    \end{problem}

    \begin{solution}
      A quadratic always has 1 turning point, since the derivative
      of a quadratic function is a linear function, which always
      changes sign once. A cubic function can have 0 or 2 turning points,
      but not 1. The derivative of a cubic is a quadratic function,
      which can have 0, 1, or 2 real roots, but if a quadratic has 1 real
      root, then it doesn't change sign at that root. This matches
      the end behavior as well. A cubic cannot have only 1 turning point,
      because then the ends would be going in the same direction.
    \end{solution}

  \item

    \note{I just want to make it explicit that cubics can have more ``shapes'' than quadratics can. (\ref{graph-cubic-with-repeated-root-followup} is a concrete example.)}

    \begin{problem}[\nstretch{0.25}]
      The graph of every quadratic can be obtained by shifting and stretching $y = x^2$. Is it true that the graph of every cubic can be obtained by shifting and stretching $y = x^3$? 
    \end{problem}

    \begin{solution}
      No, not every cubic can be graphed by shifting and stretching
      $y=x^3$. The graph of any shift and stretch of $y=x^3$ will
      only ever have one $x$-intercept, but there are cubics
      that have 2 and 3 $x$-intercepts.
    \end{solution}

  \end{enumerate}

\item  Suppose $f(x)$ is a degree $n$ polynomial.
  \begin{enumerate}

  \item
    \begin{problem}[\vfill]
      What can you say about the end behavior of $f$? (That is, what can you say about $\displaystyle \lim_{x \to \infty} f(x)$ and $\displaystyle \lim_{x \to -\infty} f(x)$?)
    \end{problem}

    \begin{solution}
      If $n$ is even, then the ends go in the same direction,
       $\lim\limits_{x\to-\infty} f(x) = \lim\limits_{x \to \infty} f(x)=
      \pm \infty$. If $n$ is odd, then the ends go in opposite directions,
      $\lim\limits_{x\to\infty} f(x)
      = \pm \infty$ and $\lim\limits_{x\to-\infty} f(x) = \mp \infty$.
      
      
    \end{solution}

  \item

    \note{$f$ can have at most $n$ roots. In addition, if $n$ is odd, $f$ must have at least 1 root because of the end behavior.}

    \begin{problem}[\vfill]
      What can you say about the number of roots of $f$?
    \end{problem}

    \begin{solution}
      The number of roots is at most $n$. If $n$ is odd, then
      $f$ must have at least 1 root (because of the end behavior.)
    \end{solution}

  \item

    \note{Since $f'$ has degree $n - 1$, $f'$ has at most $n - 1$ roots. But the end behavior also tells us something: if $n$ is even, then $f$ must have an odd number of turning points; if $n$ is odd, then $f$ must have an even number of turning points.}

    \begin{problem}[\vfill]
      What can you say about the number of turning points of $f$?
    \end{problem}

    \begin{solution}
      If $f$ has degree $n$, then $f'$ has degree $n-1$, so $f'$ has
      at most $n-1$ roots, and thus $f$ has at most $n-1$ turning points.
      Also, because of the end behavior, if $n$ is even, then there must
      be an odd number of turning points, and if $n$ is odd, then there
      must be an even number of turning points.
    \end{solution}

  \end{enumerate}

\item 

  \begin{problem}[\stretch{2}]
    Sketch the graph of $f(x) = (x - 1)(3 - x)(x - 5)$.
  \end{problem}

  \begin{solution}
    \begin{center}
      \begin{tikzpicture}
        \begin{axis}[xtick={1,3,5}, ytick=\empty]
          \addplot[domain=-1:6] {(x-1)*(3-x)*(x-5)};
        \end{axis}
      \end{tikzpicture}
    \end{center}
  \end{solution}

\item 
  \begin{problem}[\vfill\newpage]
    Find a possible formula for a degree 4 polynomial with $y$-intercept $4$ and exactly three $x$-intercepts, at $x = 1, 2, 3$. How many possible answers are there?
  \end{problem}

  \begin{solution}
    If a degree 4 polynomial has exactly 3 real roots, then it
    must have the form $a(x-r_1)^2(x-r_2)(x-r_3)$. Let's suppose
    $r_1=1$, $r_2=2$, and $r_3=3$. Then the $y$-intercept is $y=a(0-1)^2(0-2)(0-3)=6a$, so $a=\frac{2}{3}$.
    Thus, one possible formula is $\frac{2}{3}(x-1)^2(x-2)(x-3)$. We
    could find other formulas by letting $2$ or $3$ be the root with
    \textit{multiplicity} 2, i.e. the root associated with the quadratic
    factor, so there are three possible answers.
  \end{solution}

\item 

  \note{In both parts, it makes sense to start by coming up with a possible $f'$ and then using that to get $f$.}

  In each part, give an example of a cubic function $f(x)$ with the stated properties; please give a formula. These aren't problems that you can solve just by looking at them! Instead, you'll need to come up with a strategy and then carry it out.

  \begin{enumerate}
  \item
    \begin{problem}[\vfill]
      $f$ has no critical points.
    \end{problem}

    \begin{solution}
      We should start with a quadratic derivative that has no real roots,
      like $f'(x)=x^2+1$. A function with this derivative could be
      $f(x)=\frac{1}{3}x^3+x$.
    \end{solution}

  \item

    \begin{problem}[\stretch{1.25}]
      $f(x)$ has a local minimum at $x = 0$ and a local maximum at $x = 2$.
    \end{problem}

    \begin{solution}
      We want the derivative to change from negative to positive
      at $x=0$ and from positive to negative at $x=2$, so we could
      use $f'(x)=-x(x-2)=-x^2+2x$. A function with this derivative
      could be $f(x)=-\frac{1}{3}x^3 +x^2$.
    \end{solution}

  \end{enumerate}

\item 

  \note{The graph has 4 turning points, so it must have degree $\geq 5$. Based on the end behavior, it must have odd degree.}

  \begin{problem}[\newpage]
    Here's the graph of a polynomial. What can you say about its degree? Be as specific as you can. 

    \begin{tikzpicture}
      \begin{axis}[enlarge y limits=false, ytick=\empty, xtick=\empty, width=2in]
        \addplot+[domain=-2.35:3.4] {(x - 3)*(x - 2)*(x - 1)*(x + 1)*(x + 2) + 20}; % 
      \end{axis}
    \end{tikzpicture}
  \end{problem}

  \begin{solution}
    The ends are going in opposite directions, so the degree
    must be odd. There are 4 turning points, so the degree
    must be at least 5.
  \end{solution}

\item 

  Here's the graph of $f(x) = x^3 + 5x^2 + 3x - 9$ again. The critical points of $f(x)$ are $x = - 3$ and $x = - \dfrac{1}{3}$. 

  \begin{minipage}[t]{1.5in}
    \vspace{-4pt}
    \begin{tikzpicture}
      \begin{axis}[xtick={1, -3}, ytick=\empty, width=2in]
        \addplot+[domain=-5:2]{(x-1)*(x+3)^2};
      \end{axis}
    \end{tikzpicture}
  \end{minipage}
  \begin{minipage}[t]{\linewidth-1.7in}
    \vspace{-8pt}
    \begin{enumerate}
    \item
      \begin{problem}
        Find a function $g(x)$ such that $g'(x) = f(x)$ and $g(0) = 9$.
      \end{problem}

      \begin{solution}
        $g(x)=\frac{1}{4}x^4 + \frac{5}{3}x^3 + \frac{3}{2}x^2 -9x + 9$.
      \end{solution}

    \end{enumerate}
  \end{minipage}

  \begin{enumerate}[start=2]

  \item
    \begin{problem}[\vfill]
      Find and classify the critical points of $g$. 
    \end{problem}

    \begin{solution}
      $g$ has critical points at $x=-3$ and $x=1$. $x=-3$ is neither
      a local min nor a local max, and $x=1$ is a local min.
    \end{solution}

  \item
    \begin{problem}[\vfill]
      Does $g$ achieve a global minimum on $(-\infty, \infty)$? If so, where (at what $x$-value), and what's the minimum value (the $y$-value)? 
    \end{problem}

    \begin{solution}
      $g$ decreases on $(-\infty, 1)$ and increases on $(1, \infty)$,
      so the global minimum occurs at $x=1$. The minimum value is
      $g(1)=\frac{41}{12}$.
    \end{solution}

  \item
    \begin{problem}[\vfill]
      Where is $g$ concave up? Where is $g$ concave down?
    \end{problem}

    \begin{solution}
      $g$ is concave up when $g'$ is increasing, so on $(-\infty, -3)$ 
      and on $(-\frac{1}{3}, \infty)$. $g$ is concave down when $g'$
      is decreasing, so on $(-3, -\frac{1}{3})$.
    \end{solution}

  \item
    \begin{problem}[\stretch{2}]
      Sketch the graph of $g$. 
    \end{problem}

   \begin{solution}
    \begin{center}
      \begin{tikzpicture}
        \begin{axis}[xtick={-3,-1/3,1}, ytick={41/12,9},
        xticklabels={$-3$,$-1/3$,$1$}, yticklabels={$41/12$, $9$}]
          \addplot[domain=-5:2.75] {0.25*x^4+(5/3)*x^3+1.5*x^2-9*x+9};
        \end{axis}
      \end{tikzpicture}
    \end{center}
  \end{solution}

  \item
    \begin{problem}[\nstretch{0.25}]
      How many roots does $g$ have? How do you know?
    \end{problem}

    \begin{solution}
      $g$ has zero real roots. We know this because
      the global minimum value is positive.
    \end{solution}

  \end{enumerate}

\end{enumerate}
\end{document}
